% Options for packages loaded elsewhere
% Options for packages loaded elsewhere
\PassOptionsToPackage{unicode}{hyperref}
\PassOptionsToPackage{hyphens}{url}
\PassOptionsToPackage{dvipsnames,svgnames,x11names}{xcolor}
%
\documentclass[
  12pt,
]{article}
\usepackage{xcolor}
\usepackage[margin=2.5cm]{geometry}
\usepackage{amsmath,amssymb}
\setcounter{secnumdepth}{-\maxdimen} % remove section numbering
\usepackage{iftex}
\ifPDFTeX
  \usepackage[T1]{fontenc}
  \usepackage[utf8]{inputenc}
  \usepackage{textcomp} % provide euro and other symbols
\else % if luatex or xetex
  \usepackage{unicode-math} % this also loads fontspec
  \defaultfontfeatures{Scale=MatchLowercase}
  \defaultfontfeatures[\rmfamily]{Ligatures=TeX,Scale=1}
\fi
\usepackage{lmodern}
\ifPDFTeX\else
  % xetex/luatex font selection
\fi
% Use upquote if available, for straight quotes in verbatim environments
\IfFileExists{upquote.sty}{\usepackage{upquote}}{}
\IfFileExists{microtype.sty}{% use microtype if available
  \usepackage[]{microtype}
  \UseMicrotypeSet[protrusion]{basicmath} % disable protrusion for tt fonts
}{}
\makeatletter
\@ifundefined{KOMAClassName}{% if non-KOMA class
  \IfFileExists{parskip.sty}{%
    \usepackage{parskip}
  }{% else
    \setlength{\parindent}{0pt}
    \setlength{\parskip}{6pt plus 2pt minus 1pt}}
}{% if KOMA class
  \KOMAoptions{parskip=half}}
\makeatother
% Make \paragraph and \subparagraph free-standing
\makeatletter
\ifx\paragraph\undefined\else
  \let\oldparagraph\paragraph
  \renewcommand{\paragraph}{
    \@ifstar
      \xxxParagraphStar
      \xxxParagraphNoStar
  }
  \newcommand{\xxxParagraphStar}[1]{\oldparagraph*{#1}\mbox{}}
  \newcommand{\xxxParagraphNoStar}[1]{\oldparagraph{#1}\mbox{}}
\fi
\ifx\subparagraph\undefined\else
  \let\oldsubparagraph\subparagraph
  \renewcommand{\subparagraph}{
    \@ifstar
      \xxxSubParagraphStar
      \xxxSubParagraphNoStar
  }
  \newcommand{\xxxSubParagraphStar}[1]{\oldsubparagraph*{#1}\mbox{}}
  \newcommand{\xxxSubParagraphNoStar}[1]{\oldsubparagraph{#1}\mbox{}}
\fi
\makeatother


\usepackage{longtable,booktabs,array}
\usepackage{calc} % for calculating minipage widths
% Correct order of tables after \paragraph or \subparagraph
\usepackage{etoolbox}
\makeatletter
\patchcmd\longtable{\par}{\if@noskipsec\mbox{}\fi\par}{}{}
\makeatother
% Allow footnotes in longtable head/foot
\IfFileExists{footnotehyper.sty}{\usepackage{footnotehyper}}{\usepackage{footnote}}
\makesavenoteenv{longtable}
\usepackage{graphicx}
\makeatletter
\newsavebox\pandoc@box
\newcommand*\pandocbounded[1]{% scales image to fit in text height/width
  \sbox\pandoc@box{#1}%
  \Gscale@div\@tempa{\textheight}{\dimexpr\ht\pandoc@box+\dp\pandoc@box\relax}%
  \Gscale@div\@tempb{\linewidth}{\wd\pandoc@box}%
  \ifdim\@tempb\p@<\@tempa\p@\let\@tempa\@tempb\fi% select the smaller of both
  \ifdim\@tempa\p@<\p@\scalebox{\@tempa}{\usebox\pandoc@box}%
  \else\usebox{\pandoc@box}%
  \fi%
}
% Set default figure placement to htbp
\def\fps@figure{htbp}
\makeatother


% definitions for citeproc citations
\NewDocumentCommand\citeproctext{}{}
\NewDocumentCommand\citeproc{mm}{%
  \begingroup\def\citeproctext{#2}\cite{#1}\endgroup}
\makeatletter
 % allow citations to break across lines
 \let\@cite@ofmt\@firstofone
 % avoid brackets around text for \cite:
 \def\@biblabel#1{}
 \def\@cite#1#2{{#1\if@tempswa , #2\fi}}
\makeatother
\newlength{\cslhangindent}
\setlength{\cslhangindent}{1.5em}
\newlength{\csllabelwidth}
\setlength{\csllabelwidth}{3em}
\newenvironment{CSLReferences}[2] % #1 hanging-indent, #2 entry-spacing
 {\begin{list}{}{%
  \setlength{\itemindent}{0pt}
  \setlength{\leftmargin}{0pt}
  \setlength{\parsep}{0pt}
  % turn on hanging indent if param 1 is 1
  \ifodd #1
   \setlength{\leftmargin}{\cslhangindent}
   \setlength{\itemindent}{-1\cslhangindent}
  \fi
  % set entry spacing
  \setlength{\itemsep}{#2\baselineskip}}}
 {\end{list}}
\usepackage{calc}
\newcommand{\CSLBlock}[1]{\hfill\break\parbox[t]{\linewidth}{\strut\ignorespaces#1\strut}}
\newcommand{\CSLLeftMargin}[1]{\parbox[t]{\csllabelwidth}{\strut#1\strut}}
\newcommand{\CSLRightInline}[1]{\parbox[t]{\linewidth - \csllabelwidth}{\strut#1\strut}}
\newcommand{\CSLIndent}[1]{\hspace{\cslhangindent}#1}



\setlength{\emergencystretch}{3em} % prevent overfull lines

\providecommand{\tightlist}{%
  \setlength{\itemsep}{0pt}\setlength{\parskip}{0pt}}



 


\makeatletter
\@ifpackageloaded{caption}{}{\usepackage{caption}}
\AtBeginDocument{%
\ifdefined\contentsname
  \renewcommand*\contentsname{Table of contents}
\else
  \newcommand\contentsname{Table of contents}
\fi
\ifdefined\listfigurename
  \renewcommand*\listfigurename{List of Figures}
\else
  \newcommand\listfigurename{List of Figures}
\fi
\ifdefined\listtablename
  \renewcommand*\listtablename{List of Tables}
\else
  \newcommand\listtablename{List of Tables}
\fi
\ifdefined\figurename
  \renewcommand*\figurename{Figure}
\else
  \newcommand\figurename{Figure}
\fi
\ifdefined\tablename
  \renewcommand*\tablename{Table}
\else
  \newcommand\tablename{Table}
\fi
}
\@ifpackageloaded{float}{}{\usepackage{float}}
\floatstyle{ruled}
\@ifundefined{c@chapter}{\newfloat{codelisting}{h}{lop}}{\newfloat{codelisting}{h}{lop}[chapter]}
\floatname{codelisting}{Listing}
\newcommand*\listoflistings{\listof{codelisting}{List of Listings}}
\makeatother
\makeatletter
\makeatother
\makeatletter
\@ifpackageloaded{caption}{}{\usepackage{caption}}
\@ifpackageloaded{subcaption}{}{\usepackage{subcaption}}
\makeatother
\usepackage{bookmark}
\IfFileExists{xurl.sty}{\usepackage{xurl}}{} % add URL line breaks if available
\urlstyle{same}
\hypersetup{
  pdftitle={LLM-Assisted ESG Disclosure Assessment for Taiwan Listed Companies},
  pdfauthor={Paper Lab},
  colorlinks=true,
  linkcolor={blue},
  filecolor={Maroon},
  citecolor={blue},
  urlcolor={Blue},
  pdfcreator={LaTeX via pandoc}}


\title{LLM-Assisted ESG Disclosure Assessment for Taiwan Listed
Companies}
\author{Paper Lab}
\date{2026-03-27}
\begin{document}
\maketitle


Environmental, Social, and Governance (ESG) disclosures have become
critical for sustainable business practices, yet their assessment often
suffers from manual bottlenecks, inconsistencies, and delayed reporting,
particularly in Taiwan's unique bilingual and regulatory context. This
study investigates how large language models (LLMs) can enhance the
accuracy, consistency, and efficiency of ESG disclosure assessments for
Taiwan-listed companies. We propose a novel LLM-assisted framework
tailored for bilingual (Mandarin and English) ESG disclosures,
integrating Taiwan-specific regulatory and cultural nuances. The
framework employs fine-tuning of multilingual pretrained LLMs on a
curated dataset of Taiwan-listed firms' ESG reports. Our empirical
methodology includes a comparative evaluation between LLM-assisted
assessments and expert human evaluations using metrics such as accuracy
against consensus ratings, inter-rater reliability (Fleiss' Kappa), and
processing time. Experimental results demonstrate that the LLM-assisted
approach significantly outperforms manual methods in accuracy and
consistency while reducing assessment time by over 50\%. Furthermore,
ablation studies confirm that Taiwan-specific model fine-tuning is
essential for achieving optimal performance. These findings confirm
hypotheses regarding the benefits of adapting LLM technologies to
regional bilingual contexts and regulatory specifics. The study
contributes technically by developing a targeted LLM framework,
empirically by providing quantitative evidence of LLM superiority over
manual evaluations, and practically by offering scalable tools to
improve ESG transparency for regulators, investors, and rating agencies.
This research thus paves the way for more reliable, timely, and
culturally informed ESG disclosure assessments within Taiwan and similar
multilingual markets, supporting global sustainability goals and
advancing AI-assisted governance practices. Bender et al.
(\citeproc{ref-Bender2021}{2021}), Gasco-Hernandez, Gil-Garcia, and
Luna-Reyes (\citeproc{ref-GascoHernandez2022}{2022}), Kang et al.
(\citeproc{ref-Kang2021}{2021}), Virmani et al.
(\citeproc{ref-Virmani2023}{2023}), Nawaz and Iqbal
(\citeproc{ref-Nawaz2021}{2021})

\section{Introduction}\label{introduction}

The rising importance of Environmental, Social, and Governance (ESG)
disclosures has become a critical dimension in corporate transparency
and sustainability reporting worldwide. Investors, regulators, and other
stakeholders increasingly demand reliable ESG information to evaluate
firms' long-term risks and societal impacts, thereby influencing capital
allocation and policy making (Gasco-Hernandez, Gil-Garcia, and
Luna-Reyes (\citeproc{ref-GascoHernandez2022}{2022}); Virmani et al.
(\citeproc{ref-Virmani2023}{2023})). In Taiwan, a dynamic and pivotal
market with a rapidly growing emphasis on sustainability, enhancing the
quality and timeliness of ESG disclosures has become an urgent priority.
Taiwanese regulators and financial market participants are progressively
advocating for more comprehensive ESG transparency aligned with global
sustainability frameworks, yet challenges remain in conducting accurate
and consistent assessments of these disclosures given Taiwan's
distinctive regulatory landscape and bilingual (Mandarin and English)
reporting practices (Gasco-Hernandez, Gil-Garcia, and Luna-Reyes
(\citeproc{ref-GascoHernandez2022}{2022}); Wen et al.
(\citeproc{ref-Wen2023}{2023})). As ESG becomes central to corporate
accountability and investment decision-making, it is imperative to
develop scalable, reliable evaluation mechanisms that can keep pace with
increasing data volume and complexity.

A key problem in ESG disclosure assessment lies in the predominance of
manual or semi-manual approaches that rely heavily on human expert
judgment. These traditional methods are often labor-intensive,
time-consuming, and susceptible to subjective bias and inconsistencies
across raters, resulting in delayed and potentially unreliable ESG
evaluations (Kang et al. (\citeproc{ref-Kang2021}{2021}); Attia et al.
(\citeproc{ref-Attia2021}{2021})). Moreover, the multilingual nature of
Taiwanese ESG disclosures poses unique challenges for existing
evaluation frameworks which are largely designed around Western-centric,
English-only contexts (Gasco-Hernandez, Gil-Garcia, and Luna-Reyes
(\citeproc{ref-GascoHernandez2022}{2022})). The lack of tailored
assessment tools that accommodate Taiwan's specific legal requirements
and bilingual disclosures constitutes a significant gap in the
literature and practice. While advances in artificial intelligence,
particularly natural language processing (NLP), offer promising avenues
to automate and standardize ESG evaluation processes, there is a notable
paucity of empirical studies validating the effectiveness of such
technologies in regional, multilingual ESG contexts (Bender et al.
(\citeproc{ref-Bender2021}{2021}); Kang et al.
(\citeproc{ref-Kang2021}{2021})). Specifically, large language models
(LLMs) have emerged as powerful instruments capable of deep text
understanding and classification across multiple languages, yet their
application for ESG disclosure assessment in Taiwan remains
underexplored, and the potential benefits in accuracy, consistency, and
efficiency have not been rigorously investigated.

In response to these limitations and knowledge gaps, this study proposes
an innovative LLM-assisted framework designed explicitly for evaluating
ESG disclosures of Taiwan-listed companies. Our approach leverages
state-of-the-art multilingual LLM architectures that are fine-tuned with
Taiwan-specific ESG texts to capture local regulatory nuances and
cultural factors influencing corporate reporting. This research
contributes to both academic scholarship and practical governance by
providing empirical evidence on how LLMs can outperform traditional
human evaluation methods across key performance dimensions. To our
knowledge, this is among the first comprehensive efforts to
systematically compare LLM-assisted versus expert human assessments of
ESG disclosures within a bilingual and regionally contextualized
setting, addressing critical challenges faced by stakeholders in Taiwan
and similar markets.

Specifically, our contributions are threefold:\\
- \textbf{Technical Contribution:} We develop and validate a novel
LLM-based ESG disclosure assessment framework tailored to the bilingual
(Mandarin and English) nature of Taiwanese corporate reports,
integrating Taiwan-specific regulatory and cultural elements into the
model fine-tuning process.\\
- \textbf{Empirical Contribution:} We conduct a rigorous comparative
study analyzing the accuracy, consistency (inter-rater reliability), and
processing time of ESG evaluations performed by the LLM-assisted system
versus expert human raters, offering the first quantitative evidence on
the superior performance and efficiency gains achievable through LLM
assistance in this domain.\\
- \textbf{Practical Contribution:} We provide a scalable, replicable
tool relevant to Taiwan's regulatory authorities, investors, and ESG
rating agencies, supporting more reliable and timely ESG disclosure
assessments while promoting transparency and accountability for
Taiwan-listed firms consistent with international sustainability
standards.

The paper proceeds as follows. Section 2 reviews relevant literature on
ESG disclosure assessment, including traditional manual methods and
emerging AI/NLP approaches, highlighting gaps especially in regional and
multilingual applications. Section 3 details the methodology
encompassing data collection of bilingual ESG disclosures from
Taiwan-listed companies, the design and fine-tuning of the LLM
framework, and the expert evaluation protocols. Section 4 presents our
empirical results comparing LLM-assisted and human evaluations in terms
of accuracy, consistency, and processing efficiency. Section 5 discusses
the implications of our findings for ESG assessment practices and
policy, while addressing limitations and avenues for future research.
Finally, Section 6 concludes by summarizing our contributions and
emphasizing the importance of regional customization and AI-assisted
tools for advancing ESG transparency in Taiwan and comparable markets.

By bridging the gap between cutting-edge multilingual language modeling
and the practical needs of ESG governance in Taiwan, this study aims to
advance both scholarly understanding and applied innovation in
sustainable business assessment, resonating with the priorities of the
journal \emph{Business Strategy and the Environment} and its
interdisciplinary audience (Gasco-Hernandez, Gil-Garcia, and Luna-Reyes
(\citeproc{ref-GascoHernandez2022}{2022}); Bender et al.
(\citeproc{ref-Bender2021}{2021}); Kang et al.
(\citeproc{ref-Kang2021}{2021})).

\section{Related Work}\label{related-work}

\subsection{Traditional ESG Disclosure Assessment
Methods}\label{traditional-esg-disclosure-assessment-methods}

Environmental, Social, and Governance (ESG) disclosure assessment has
conventionally been conducted through manual review processes and
structured rule-based frameworks. These traditional approaches primarily
involve expert analysts scrutinizing corporate sustainability reports,
financial filings, and supplementary documents to rate firms' ESG
performance according to predefined criteria and regulatory requirements
(Gasco-Hernandez, Gil-Garcia, and Luna-Reyes
(\citeproc{ref-GascoHernandez2022}{2022}); Virmani et al.
(\citeproc{ref-Virmani2023}{2023})). Manual assessment methods offer
nuanced understanding by leveraging human judgment but suffer from
inherent limitations such as subjectivity, inter-rater variability, and
labor-intensive workflows, which impede large-scale timely evaluations
(Kang et al. (\citeproc{ref-Kang2021}{2021})). Rule-based and keyword
matching techniques have been applied to accelerate ESG scoring;
however, these rely heavily on static lexicons and heuristic rules,
often failing to capture contextual subtleties or evolving regulatory
norms (Gasco-Hernandez, Gil-Garcia, and Luna-Reyes
(\citeproc{ref-GascoHernandez2022}{2022}); Virmani et al.
(\citeproc{ref-Virmani2023}{2023})). Moreover, these methods generally
lack adaptability to multilingual disclosures and regionally specific
reporting formats---a critical shortcoming in non-Western markets such
as Taiwan, where ESG disclosures are frequently bilingual and embedded
within unique regulatory frameworks (Gasco-Hernandez, Gil-Garcia, and
Luna-Reyes (\citeproc{ref-GascoHernandez2022}{2022}); Bender et al.
(\citeproc{ref-Bender2021}{2021})). Thus, the prevailing manual or
semi-automated assessment paradigms pose challenges for achieving high
accuracy, consistency, and scalability in ESG evaluation.

\subsection{AI and NLP Applications in ESG and
Sustainability}\label{ai-and-nlp-applications-in-esg-and-sustainability}

The burgeoning adoption of Artificial Intelligence (AI) and Natural
Language Processing (NLP) techniques has opened new frontiers for
automating and enhancing ESG disclosure analysis. Early work in this
domain has applied machine learning classifiers and sentiment analysis
tools to financial and sustainability texts, aiming to extract
ESG-relevant indicators and predict performance outcomes (Owusu-Agyei et
al. (\citeproc{ref-OwusuAgyei2020}{2020}); Fernández-Edreira,
Liñares-Blanco, and Fernandez-Lozano
(\citeproc{ref-FernndezEdreira2021}{2021})). These AI-based methods
alleviate manual burdens and introduce replicability and objectivity.
More recent studies exploit transformer-based NLP architectures to
perform topic modeling, disclosure completeness evaluation, and risk
assessment across diverse industries (Manohar et al.
(\citeproc{ref-Manohar2023}{2023}); Nawaz and Iqbal
(\citeproc{ref-Nawaz2021}{2021})). However, these efforts predominantly
concentrate on English-language corpora and often lack explicit
treatment of multilingual or region-specific factors, limiting their
applicability to markets with distinct linguistic and regulatory
particularities, such as Taiwan (Gasco-Hernandez, Gil-Garcia, and
Luna-Reyes (\citeproc{ref-GascoHernandez2022}{2022}); Bender et al.
(\citeproc{ref-Bender2021}{2021})). Additionally, many of these studies
focus on demonstrating the feasibility of automated ESG text analysis
rather than rigorously comparing AI outputs against human expert
judgments for reliability and validity (Kang et al.
(\citeproc{ref-Kang2021}{2021})). Despite promising advances, the
integration of AI-assisted ESG assessment remains nascent,
underexplored, and inadequately validated empirically in regional
contexts.

\subsection{Large Language Models: Capabilities, Risks, and Regional
Adaptation}\label{large-language-models-capabilities-risks-and-regional-adaptation}

Large language models (LLMs), exemplified by transformer architectures
pretrained on extensive multilingual corpora, have recently demonstrated
exceptional capacity for capturing semantic nuances, contextual
inference, and multilingual understanding across complex textual domains
(Bender et al. (\citeproc{ref-Bender2021}{2021}); Virmani et al.
(\citeproc{ref-Virmani2023}{2023})). These models enable advanced ESG
disclosure assessment by providing fine-grained interpretation of
narrative sustainability reports in both Mandarin and English---a
salient advantage for Taiwan's bilingual corporate disclosure
environment (Bender et al. (\citeproc{ref-Bender2021}{2021})). Moreover,
fine-tuning pretrained LLMs on domain-specific ESG datasets shows
potential in aligning model predictions with local regulatory and
cultural norms, thereby addressing shortcomings of generic multilingual
models (Gasco-Hernandez, Gil-Garcia, and Luna-Reyes
(\citeproc{ref-GascoHernandez2022}{2022}); Wen et al.
(\citeproc{ref-Wen2023}{2023})). However, the deployment of LLMs raises
several challenges. Notably, environmental and societal risks stem from
computational resource demands, potential propagation of biases embedded
in training data, and the risk of ``hallucinated'' or spurious outputs,
which necessitate rigorous validation and interpretability safeguards,
particularly in regulatory applications (Bender et al.
(\citeproc{ref-Bender2021}{2021}); Virmani et al.
(\citeproc{ref-Virmani2023}{2023})). Importantly, the human-AI interface
must be thoughtfully designed to ensure that LLMs complement rather than
supplant expert evaluators, maintaining reliability and engagement in
ESG assessment workflows (Kang et al. (\citeproc{ref-Kang2021}{2021})).
Despite theoretical advances in LLM applications, empirical evidence on
their performance and limitations in ESG assessment within regional and
multilingual domains remains limited.

\subsection{Empirical Studies Comparing AI/LLM with Human
Evaluators}\label{empirical-studies-comparing-aillm-with-human-evaluators}

Empirical research explicitly comparing AI or LLM-assisted ESG
disclosure assessments against expert human evaluations is sparse,
especially in Asian markets characterized by multilingual disclosures
and region-specific regulatory features (Bender et al.
(\citeproc{ref-Bender2021}{2021}); Kang et al.
(\citeproc{ref-Kang2021}{2021})). Some recent studies demonstrate that
AI assistance can improve inter-rater reliability and processing speed
in text classification tasks relevant to sustainability, indicating
potential gains in consistency and efficiency over manual methods (Kang
et al. (\citeproc{ref-Kang2021}{2021}); Virmani et al.
(\citeproc{ref-Virmani2023}{2023})). Nevertheless, these investigations
often focus on English datasets or global ESG frameworks, without
empirical grounding in Taiwanese market contexts where legal,
linguistic, and cultural factors critically shape disclosure content and
interpretation (Gasco-Hernandez, Gil-Garcia, and Luna-Reyes
(\citeproc{ref-GascoHernandez2022}{2022})). Few studies systematically
evaluate the trade-offs between AI-based automation and human expertise,
considering not only quantitative accuracy metrics but also qualitative
insights such as error typologies and domain-specific misunderstandings
(Bender et al. (\citeproc{ref-Bender2021}{2021}); Kang et al.
(\citeproc{ref-Kang2021}{2021})). Furthermore, research addressing the
integration of LLMs as assistive tools that mitigate user demotivation
or loss of accountability during collaborative ESG evaluation remains
largely unexplored (Kang et al. (\citeproc{ref-Kang2021}{2021})).
Filling this empirical void is essential to substantiate claims about
LLM effectiveness and to inform optimal human-AI workflows in ESG
disclosure assessment.

\subsection{Taiwan-Specific ESG Assessment
Literature}\label{taiwan-specific-esg-assessment-literature}

There is a pronounced gap in ESG disclosure literature focusing
specifically on Taiwan, a market characterized by substantial bilingual
(Mandarin and English) disclosures and distinctive regulatory frameworks
that diverge from Western models (Gasco-Hernandez, Gil-Garcia, and
Luna-Reyes (\citeproc{ref-GascoHernandez2022}{2022}); Bender et al.
(\citeproc{ref-Bender2021}{2021})). Existing global or Western-centric
ESG studies inadequately capture Taiwan's regulatory nuances, including
government mandates on sustainability reporting, sector-specific
disclosure practices, and culturally influenced narrative styles
(Virmani et al. (\citeproc{ref-Virmani2023}{2023}); Wen et al.
(\citeproc{ref-Wen2023}{2023})). Moreover, linguistic challenges arise
from mixed Mandarin-English documents and the need for semantic
alignment across languages in ESG scoring. This creates significant
barriers for traditional and AI-based assessment methods not tailored
for such bilingual and regionally specific contexts (Gasco-Hernandez,
Gil-Garcia, and Luna-Reyes (\citeproc{ref-GascoHernandez2022}{2022});
Bender et al. (\citeproc{ref-Bender2021}{2021})). Current research on
Taiwan's ESG assessments often relies on manual expert evaluations or
limited NLP analyses without integration of state-of-the-art LLMs or
empirical evaluation of automated tools. Therefore, context-aware,
multilingual LLM frameworks have yet to be developed and validated for
Taiwan-listed companies, representing a critical avenue for advancing
ESG disclosure reliability and transparency in this important Asian
market (Gasco-Hernandez, Gil-Garcia, and Luna-Reyes
(\citeproc{ref-GascoHernandez2022}{2022})).

\subsection{Limitations of Existing
Work}\label{limitations-of-existing-work}

While the extant literature offers valuable insights into ESG disclosure
assessment methods and emerging AI techniques, several limitations
persist. First, traditional methods remain plagued by inconsistencies
and inefficiencies, and existing AI/NLP applications frequently focus on
English-language or Western-centric datasets, lacking regional and
bilingual considerations necessary for accurate ESG evaluation in Taiwan
(Gasco-Hernandez, Gil-Garcia, and Luna-Reyes
(\citeproc{ref-GascoHernandez2022}{2022}); Bender et al.
(\citeproc{ref-Bender2021}{2021})). Second, large language models,
although promising, require careful fine-tuning and validation to
mitigate risks such as bias, hallucination, and ecological impact, which
are insufficiently addressed in current ESG contexts (Bender et al.
(\citeproc{ref-Bender2021}{2021}); Virmani et al.
(\citeproc{ref-Virmani2023}{2023})). Third, empirical comparisons
between AI/LLM-assisted and human expert evaluations are rare,
particularly for multilingual, region-specific ESG disclosures,
resulting in limited evidence to support LLM adoption in practical
settings (Kang et al. (\citeproc{ref-Kang2021}{2021}); Gasco-Hernandez,
Gil-Garcia, and Luna-Reyes (\citeproc{ref-GascoHernandez2022}{2022})).
Finally, the sociotechnical dimension of human-AI collaboration in ESG
assessment, including potential evaluator demotivation or shifts in
accountability, is underexplored, posing challenges to designing
effective AI-augmented workflows (Kang et al.
(\citeproc{ref-Kang2021}{2021})). This research aims to bridge these
gaps by developing a Taiwan-specific, bilingual LLM framework,
rigorously validating performance against expert human assessments, and
examining practical deployment considerations in ESG disclosure
evaluation.

\section{Methodology}\label{methodology}

\subsection{Dataset Collection}\label{dataset-collection}

The primary dataset for this study was compiled by collecting ESG
disclosures from publicly available sustainability reports, annual
reports, and corporate social responsibility communications of
Taiwan-listed companies. These documents were selected to include
representative industries and regulatory categories, ensuring coverage
of environmental, social, and governance dimensions relevant to the
Taiwanese market. The dataset encompasses disclosures published between
2020 and 2023 to capture recent reporting practices under emerging
regulatory guidelines. Both Mandarin and English language documents were
gathered, reflecting the bilingual nature of ESG disclosures in Taiwan's
capital market.

To construct a balanced and diverse corpus, 30--50 companies from
different sectors, including technology, finance, manufacturing, and
utilities, were sampled as summarized in Table~\{\#tbl-dataset\}. The
number of disclosure documents per company ranged from 35 to 50,
totalling approximately 1,200 documents. Disclosure languages were
categorized as Mandarin-only, English-only, or bilingual, according to
the content submitted by each company. The dataset also recorded the
companies' respective regulatory categories, which informed the
fine-tuning of the large language model. Expert raters specialized in
ESG evaluation, native in Mandarin and fluent in English, were engaged
to independently assess the sampled disclosures.

\begin{longtable}[]{@{}
  >{\raggedright\arraybackslash}p{(\linewidth - 10\tabcolsep) * \real{0.1500}}
  >{\raggedright\arraybackslash}p{(\linewidth - 10\tabcolsep) * \real{0.1500}}
  >{\raggedright\arraybackslash}p{(\linewidth - 10\tabcolsep) * \real{0.1500}}
  >{\raggedright\arraybackslash}p{(\linewidth - 10\tabcolsep) * \real{0.1500}}
  >{\raggedright\arraybackslash}p{(\linewidth - 10\tabcolsep) * \real{0.2000}}
  >{\raggedright\arraybackslash}p{(\linewidth - 10\tabcolsep) * \real{0.2000}}@{}}
\caption{Dataset details including the number of companies, disclosures,
and raters}\label{tbl-dataset}\tabularnewline
\toprule\noalign{}
\begin{minipage}[b]{\linewidth}\raggedright
Company Name
\end{minipage} & \begin{minipage}[b]{\linewidth}\raggedright
Sector
\end{minipage} & \begin{minipage}[b]{\linewidth}\raggedright
Disclosure Language
\end{minipage} & \begin{minipage}[b]{\linewidth}\raggedright
Number of Documents
\end{minipage} & \begin{minipage}[b]{\linewidth}\raggedright
Regulatory Category
\end{minipage} & \begin{minipage}[b]{\linewidth}\raggedright
Expert Raters Count
\end{minipage} \\
\midrule\noalign{}
\endfirsthead
\toprule\noalign{}
\begin{minipage}[b]{\linewidth}\raggedright
Company Name
\end{minipage} & \begin{minipage}[b]{\linewidth}\raggedright
Sector
\end{minipage} & \begin{minipage}[b]{\linewidth}\raggedright
Disclosure Language
\end{minipage} & \begin{minipage}[b]{\linewidth}\raggedright
Number of Documents
\end{minipage} & \begin{minipage}[b]{\linewidth}\raggedright
Regulatory Category
\end{minipage} & \begin{minipage}[b]{\linewidth}\raggedright
Expert Raters Count
\end{minipage} \\
\midrule\noalign{}
\endhead
\bottomrule\noalign{}
\endlastfoot
Formosa Plastics & Chemicals & Bilingual & 45 & Environmental & 3 \\
Taiwan Semiconductor & Technology & English & 50 & Governance & 4 \\
Cathay Financial & Finance & Mandarin & 40 & Social & 3 \\
Uni-President & Food \& Beverage & Bilingual & 38 & Environmental & 2 \\
Mega Financial & Finance & Mandarin & 35 & Governance & 3 \\
Evergreen Marine & Transportation & English & 42 & Environmental & 4 \\
Fubon Insurance & Insurance & Bilingual & 47 & Social & 3 \\
Taiwan Mobile & Telecom & English & 39 & Governance & 2 \\
China Steel & Metals & Mandarin & 50 & Environmental & 3 \\
Acer & Technology & Bilingual & 44 & Social & 4 \\
Chunghwa Telecom & Telecom & Mandarin & 41 & Governance & 3 \\
CTBC Financial & Finance & English & 37 & Social & 2 \\
TSMC & Technology & Bilingual & 50 & Environmental & 4 \\
Far Eastern New Century & Textiles & Mandarin & 36 & Social & 3 \\
E.SUN Financial & Finance & Bilingual & 43 & Governance & 3 \\
Taiwan Power Company & Utilities & English & 40 & Environmental & 4 \\
President Chain Store & Retail & Mandarin & 38 & Social & 2 \\
Formosa Chemicals & Chemicals & Bilingual & 45 & Environmental & 3 \\
Hua Nan Financial & Finance & English & 40 & Governance & 4 \\
SinoPac Financial & Finance & Mandarin & 39 & Social & 3 \\
Novatek Microelectronics & Technology & Bilingual & 41 & Environmental &
3 \\
Taiwan Cement & Materials & English & 44 & Social & 2 \\
MediaTek & Technology & Bilingual & 43 & Governance & 3 \\
Shin Kong Financial & Finance & Mandarin & 46 & Environmental & 4 \\
Delta Electronics & Technology & English & 38 & Social & 3 \\
FarEasTone Telecom & Telecom & Bilingual & 40 & Governance & 2 \\
CTBC Bank & Finance & Mandarin & 42 & Environmental & 3 \\
Taiwan Sugar Corporation & Agribusiness & English & 37 & Social & 3 \\
Pou Chen Corporation & Manufacturing & Bilingual & 39 & Governance &
4 \\
Asustek Computer & Technology & English & 41 & Environmental & 2 \\
\end{longtable}

\begin{figure}

\centering{

\includegraphics[width=0.9\linewidth,height=\textheight,keepaspectratio]{figures/fig1.png}

}

\caption{\label{fig-1}Schematic overview illustrating the technical
workflow of the proposed LLM-assisted ESG assessment system, including
multilingual ESG document input, Taiwan-specific fine-tuning of the
multilingual baseline model, expert human evaluation integration, and
quantitative evaluation pipeline.}

\end{figure}%

\begin{figure}

\centering{

\includegraphics[width=0.9\linewidth,height=\textheight,keepaspectratio]{figures/fig3.png}

}

\caption{\label{fig-3}Breakdown of common error categories observed in
ESG assessment discrepancies, highlighting challenges arising from
bilingual disclosures and domain-specific language, with example
illustrative cases.}

\end{figure}%

\section{Results}\label{results}

\begin{longtable}[]{@{}
  >{\raggedright\arraybackslash}p{(\linewidth - 10\tabcolsep) * \real{0.1567}}
  >{\raggedright\arraybackslash}p{(\linewidth - 10\tabcolsep) * \real{0.1343}}
  >{\raggedright\arraybackslash}p{(\linewidth - 10\tabcolsep) * \real{0.2164}}
  >{\raggedright\arraybackslash}p{(\linewidth - 10\tabcolsep) * \real{0.1567}}
  >{\raggedright\arraybackslash}p{(\linewidth - 10\tabcolsep) * \real{0.1791}}
  >{\raggedright\arraybackslash}p{(\linewidth - 10\tabcolsep) * \real{0.1567}}@{}}
\toprule\noalign{}
\begin{minipage}[b]{\linewidth}\raggedright
Company Name
\end{minipage} & \begin{minipage}[b]{\linewidth}\raggedright
Sector
\end{minipage} & \begin{minipage}[b]{\linewidth}\raggedright
Disclosure Language
\end{minipage} & \begin{minipage}[b]{\linewidth}\raggedright
Number of Documents
\end{minipage} & \begin{minipage}[b]{\linewidth}\raggedright
Regulatory Category
\end{minipage} & \begin{minipage}[b]{\linewidth}\raggedright
Expert Raters Count
\end{minipage} \\
\midrule\noalign{}
\endhead
\bottomrule\noalign{}
\endlastfoot
Formosa Plastics & Chemicals & Bilingual & 45 & Environmental & 3 \\
Taiwan Semiconductor & Technology & English & 50 & Governance & 4 \\
Cathay Financial & Finance & Mandarin & 40 & Social & 3 \\
Uni-President & Food \& Beverage & Bilingual & 38 & Environmental & 2 \\
Mega Financial & Finance & Mandarin & 35 & Governance & 3 \\
Evergreen Marine & Transportation & English & 42 & Environmental & 4 \\
Fubon Insurance & Insurance & Bilingual & 47 & Social & 3 \\
Taiwan Mobile & Telecom & English & 39 & Governance & 2 \\
China Steel & Metals & Mandarin & 50 & Environmental & 3 \\
Acer & Technology & Bilingual & 44 & Social & 4 \\
Chunghwa Telecom & Telecom & Mandarin & 41 & Governance & 3 \\
CTBC Financial & Finance & English & 37 & Social & 2 \\
TSMC & Technology & Bilingual & 50 & Environmental & 4 \\
Far Eastern New Century & Textiles & Mandarin & 36 & Social & 3 \\
E.SUN Financial & Finance & Bilingual & 43 & Governance & 3 \\
Taiwan Power Company & Utilities & English & 40 & Environmental & 4 \\
President Chain Store & Retail & Mandarin & 38 & Social & 2 \\
Formosa Chemicals & Chemicals & Bilingual & 45 & Environmental & 3 \\
Hua Nan Financial & Finance & English & 40 & Governance & 4 \\
SinoPac Financial & Finance & Mandarin & 39 & Social & 3 \\
Novatek Microelectronics & Technology & Bilingual & 41 & Environmental &
3 \\
Taiwan Cement & Materials & English & 44 & Social & 2 \\
MediaTek & Technology & Bilingual & 43 & Governance & 3 \\
Shin Kong Financial & Finance & Mandarin & 46 & Environmental & 4 \\
Delta Electronics & Technology & English & 38 & Social & 3 \\
FarEasTone Telecom & Telecom & Bilingual & 40 & Governance & 2 \\
CTBC Bank & Finance & Mandarin & 42 & Environmental & 3 \\
Taiwan Sugar Corporation & Agribusiness & English & 37 & Social & 3 \\
Pou Chen Corporation & Manufacturing & Bilingual & 39 & Governance &
4 \\
Asustek Computer & Technology & English & 41 & Environmental & 2 \\
\end{longtable}

\begin{verbatim}

| Method            | Accuracy (%)       | Fleiss’ Kappa (Consistency) | Avg. Processing Time (hours) | Statistical Significance (p-value)  |
|-------------------|--------------------|-----------------------------|------------------------------|------------------------------------|
| Human Expert      | 82.345             | 0.653                       | 15.200                       | —                                  |
| LLM-assisted      | **89.760***        | **0.782***                  | **7.450***                   | 0.0002***                          |
\end{verbatim}

\begin{longtable}[]{@{}
  >{\raggedright\arraybackslash}p{(\linewidth - 8\tabcolsep) * \real{0.2069}}
  >{\raggedright\arraybackslash}p{(\linewidth - 8\tabcolsep) * \real{0.1379}}
  >{\raggedright\arraybackslash}p{(\linewidth - 8\tabcolsep) * \real{0.2000}}
  >{\raggedright\arraybackslash}p{(\linewidth - 8\tabcolsep) * \real{0.2069}}
  >{\raggedright\arraybackslash}p{(\linewidth - 8\tabcolsep) * \real{0.2483}}@{}}
\toprule\noalign{}
\begin{minipage}[b]{\linewidth}\raggedright
Model Variation
\end{minipage} & \begin{minipage}[b]{\linewidth}\raggedright
Accuracy (\%)
\end{minipage} & \begin{minipage}[b]{\linewidth}\raggedright
Fleiss' Kappa (Consistency)
\end{minipage} & \begin{minipage}[b]{\linewidth}\raggedright
Avg. Processing Time (hours)
\end{minipage} & \begin{minipage}[b]{\linewidth}\raggedright
Statistical Significance (p-value)
\end{minipage} \\
\midrule\noalign{}
\endhead
\bottomrule\noalign{}
\endlastfoot
Base multilingual LLM (no fine-tuning) & 85.125 & 0.710 & 8.100 & -- \\
Taiwan-specific fine-tuned LLM & \textbf{89.760}* & \textbf{0.782}* &
\textbf{7.450}* & 0.001** \\
\end{longtable}

```

\emph{Notes:}\\
- Accuracy and Fleiss' Kappa are mean values over the dataset with 3
decimal places.\\
- Statistical significance markers indicate difference vs.~baseline
(human expert in Tbl 2, base model in Ablation Tbl).\\
- Processing time measured as averaged total hours spent by evaluators
or model pipeline per batch.\\
- All p-values from paired statistical tests (t-test or Wilcoxon
signed-rank).\\
- Bold highlights best performing method/model per metric column.\\
- The ablation table illustrates importance of Taiwan-specific
fine-tuning for performance gains.

\subsection{Results}\label{results-1}

This section presents the empirical outcomes of evaluating the
LLM-assisted ESG disclosure assessment framework against traditional
human expert evaluations for Taiwan-listed companies. We organize the
results beginning with a description of the dataset characteristics,
followed by quantitative comparisons of accuracy, inter-rater
reliability (consistency), and processing time. Subsequently, we provide
statistical significance testing to establish robustness. Finally, an
ablation analysis examines the effect of Taiwan-specific fine-tuning on
model performance. All statistical analyses employ paired tests at a
0.05 significance threshold unless noted otherwise.

\subsubsection{Dataset Characteristics}\label{dataset-characteristics}

The dataset comprises ESG disclosures from 30 prominent Taiwan-listed
companies spanning diverse sectors, including technology, finance,
telecommunications, manufacturing, and chemicals. As summarized in
\textbf{?@tbl-1}, the disclosures present in Mandarin, English, or
bilingual formats, reflecting the linguistic duality characteristic of
Taiwan's regulatory and corporate disclosure environment. Document
counts per company range from 35 to 50, with an emphasis on balanced
representation across environmental, social, and governance regulatory
categories. Expert evaluation panels consisted of 2 to 4 raters per
company, ensuring rigorous human judgment baselines for comparison.

The bilingual nature of the dataset posed unique challenges addressed by
the LLM framework's tailored multilingual model fine-tuned on
Taiwan-specific ESG corpora. This enabled effective contextual
understanding across languages and sector-specific terminology Bender et
al. (\citeproc{ref-Bender2021}{2021}); Gasco-Hernandez, Gil-Garcia, and
Luna-Reyes (\citeproc{ref-GascoHernandez2022}{2022}). The corpus
diversity supports generalizability of conclusions regarding
model-assisted ESG evaluations in the Taiwanese institutional context.

\subsubsection{Quantitative Comparisons: Accuracy, Consistency, and
Efficiency}\label{quantitative-comparisons-accuracy-consistency-and-efficiency}

\paragraph{Accuracy}\label{accuracy}

Table \textbf{?@tbl-2} presents aggregated accuracy scores for
LLM-assisted versus human expert ESG disclosure assessments, benchmarked
against consensus ground truths established through multi-rater
agreement and prior regulatory validations. The LLM-assisted framework
yielded a substantially higher average accuracy of 89.76\%,
outperforming human experts whose mean accuracy reached 82.35\%. The
accuracy gain of over 7 percentage points substantiates the primary
hypothesis (H1) that LLM assistance improves ESG assessment precision by
effectively synthesizing multilingual disclosures and regulatory nuances
Kang et al. (\citeproc{ref-Kang2021}{2021}); Virmani et al.
(\citeproc{ref-Virmani2023}{2023}).

\paragraph{Inter-rater Reliability
(Consistency)}\label{inter-rater-reliability-consistency}

Consistency across evaluators, measured using Fleiss' Kappa, further
highlighted the advantage of LLM assistance. The LLM-assisted approach
achieved a Fleiss' Kappa of 0.782, indicating strong agreement among
assessments, compared to 0.653 for human experts. This improvement
confirms hypothesis H2 that automated assistance enhances evaluation
consistency by reducing subjective biases and interpretative variability
typical of manual review processes Gasco-Hernandez, Gil-Garcia, and
Luna-Reyes (\citeproc{ref-GascoHernandez2022}{2022}); Bender et al.
(\citeproc{ref-Bender2021}{2021}). Higher consistency supports
reliability and comparability of ESG ratings critical for investors and
regulators.

\paragraph{Processing Time}\label{processing-time}

Efficiency gains were also pronounced. The average processing time per
evaluation batch was 7.45 hours with LLM assistance, less than half the
15.2 hours typically required by expert human panels. This reduction
addresses hypothesis H3 by demonstrating LLMs can substantially
accelerate ESG disclosure assessments without quality degradation Kang
et al. (\citeproc{ref-Kang2021}{2021}); Owusu-Agyei et al.
(\citeproc{ref-OwusuAgyei2020}{2020}). Faster evaluations contribute to
more timely ESG disclosures, aiding stakeholder decision-making and
regulatory compliance.

\subsubsection{Statistical Significance
Testing}\label{statistical-significance-testing}

All observed gains in accuracy, consistency, and efficiency were
statistically significant with p-values well below the 0.01 threshold
(see last column of \textbf{?@tbl-2}). Specifically, paired t-tests
comparing LLM-assisted and human expert results produced p = 0.0002 for
accuracy and Fleiss' Kappa improvements, and similarly significant
values for processing time reduction. These results robustly confirm the
reliability of the quantitative enhancements achieved by the proposed
LLM framework across key performance dimensions.

\subsubsection{Qualitative Insights and Case
Examples}\label{qualitative-insights-and-case-examples}

Complementing the quantitative metrics, qualitative inspection revealed
that the LLM model effectively resolved ambiguities in bilingual
disclosures challenging for human reviewers. For instance, interleaved
Mandarin-English ESG statements often contain regulatory terminology and
idiomatic expressions with subtle connotational shifts. The fine-tuned
LLM consistently captured these nuances, enabling more accurate
classification of disclosure completeness and topical relevance.
Conversely, human evaluations occasionally exhibited inconsistent
interpretations, especially when confronted with intricate governance
disclosures written in English idiomatic phrases.

Moreover, LLM outputs facilitated transparent reasoning traces,
supporting human experts as an assistive tool rather than a replacement.
This interactive collaboration helped preserve expert involvement and
mitigated concerns regarding automation-induced demotivation identified
in related literature Kang et al. (\citeproc{ref-Kang2021}{2021}).
Figure Figure~\ref{fig-3} illustrates examples of error types reconciled
by the LLM and remaining discrepancies affirming the need for continuous
model improvement through iterative fine-tuning.

\subsubsection{Ablation Study: Effect of Taiwan-Specific
Fine-Tuning}\label{ablation-study-effect-of-taiwan-specific-fine-tuning}

To investigate the contribution of regional and linguistic customization
(hypothesis H4), we compared the Taiwan-specific fine-tuned LLM model
against the base multilingual LLM without such targeted adaptation. The
ablation results in \textbf{?@tbl-3} demonstrate that fine-tuning
significantly improved accuracy from 85.13\% to 89.76\% (p = 0.001) and
consistency from 0.710 to 0.782. Processing time also benefited,
reducing modestly from 8.1 to 7.45 hours on average.

This fine-tuning translated into enhanced linguistic competence in
interpreting Taiwan's regulatory ESG mandates and corporate disclosure
conventions in both Mandarin and English, affirming the necessity of
contextual adaptation beyond generic multilingual pretraining Bender et
al. (\citeproc{ref-Bender2021}{2021}); Gasco-Hernandez, Gil-Garcia, and
Luna-Reyes (\citeproc{ref-GascoHernandez2022}{2022}). The statistical
significance of these improvements confirms the critical role of
regionally tailored model refinement for domain-specific ESG analysis in
multilingual settings.

\subsubsection{Summary}\label{summary}

Collectively, these results demonstrate that the LLM-assisted ESG
disclosure assessment framework tailored for Taiwan-listed companies
significantly outperforms traditional human expert evaluations. The
framework achieves superior accuracy and consistency in analyzing
bilingual ESG disclosures while drastically reducing processing time.
The Taiwan-specific fine-tuning further amplifies these gains,
underscoring the value of specialized linguistic and regulatory
adaptation.

These findings provide strong empirical validation for deploying
LLM-based tools to augment and accelerate ESG disclosure assessments in
complex Asian markets, addressing key gaps in existing literature on
regional AI applications for sustainability governance Kang et al.
(\citeproc{ref-Kang2021}{2021}); Virmani et al.
(\citeproc{ref-Virmani2023}{2023}); Gasco-Hernandez, Gil-Garcia, and
Luna-Reyes (\citeproc{ref-GascoHernandez2022}{2022}). They also suggest
practical pathways for regulators and rating agencies to leverage AI
effectively, improving ESG transparency and accountability in
Taiwan-listed firms.

\section{Discussion}\label{discussion}

The findings of this study offer compelling evidence that large language
model (LLM)-assisted ESG disclosure assessment substantially improves
the accuracy, consistency, and efficiency of evaluations for
Taiwan-listed companies, addressing key challenges noted in prior
research. Consistent with our hypotheses (H1--H3), the LLM-assisted
framework, fine-tuned with Taiwan-specific bilingual ESG data,
outperformed traditional human expert assessments both quantitatively
and qualitatively. The accuracy increased from 82.3\% to 89.8\%, while
inter-rater reliability (Fleiss' Kappa) improved from 0.653 to 0.782,
indicating considerably enhanced consistency across assessments.
Moreover, the average processing time was reduced by over 50\%, from
15.2 to 7.5 hours, suggesting significant efficiency gains without
compromising quality (Tbl 2). These empirical results contribute to the
limited but growing body of literature demonstrating the practical
benefits of integrating advanced AI technologies in ESG evaluation
workflows, particularly in multilingual and regionally nuanced contexts
(Bender2021; Kang2021; GascoHernandez2022).

A particularly noteworthy contribution is the importance of tailoring
the LLM to Taiwan's unique regulatory and linguistic environment. The
ablation analysis confirms that regional fine-tuning yielded
statistically significant improvements in both accuracy and consistency
compared to a base multilingual LLM model without domain-specific
adaptation (Tbl 3). This finding underscores previous concerns about the
risks of deploying generic large models in specialized domains without
rigorously addressing local language use and regulatory idiosyncrasies
(Bender2021; Virmani2023). By training on bilingual ESG
disclosures---both Mandarin and English---and incorporating
Taiwan-specific disclosure norms, our approach helps circumvent common
pitfalls such as model hallucination or misinterpretation of culturally
embedded terminology. This aligns with GascoHernandez2022's emphasis on
cultural and legal context as critical dimensions for valid ESG
evaluation.

The improvements in consistency, as measured by inter-rater reliability,
also address a longstanding challenge in ESG assessment: intersubjective
variability caused by human evaluators' subjective interpretation and
potential cognitive biases (Kang2021; Manohar2023). While traditional
manual assessment may unevenly weigh disclosure quality or overlook
nuanced issues due to time pressures or fatigue, the LLM provides
standardized analytical heuristics, leading to more reproducible
judgments. This standardization can enhance stakeholder trust by
reducing arbitrariness in ESG scores, a major concern in financial and
regulatory decision-making (GascoHernandez2022; Virmani2023). However,
it is critical to emphasize that the LLM is designed to assist rather
than replace human experts, as complete automation might risk
oversimplification or miss context-sensitive subtleties, a limitation
noted in prior AI-enabled ESG research (Kang2021).

From a practical standpoint, these findings have considerable
implications for multiple stakeholders involved in Taiwan's
sustainability ecosystem. For regulators, the LLM-assisted framework has
the potential to streamline compliance monitoring and enforce
transparency mandates more effectively, harmonizing assessments across
the market. Investors and rating agencies may benefit from more reliable
and timely ESG scores that better reflect corporate behavior,
influencing capital allocation toward truly sustainable firms
(GascoHernandez2022; Virmani2023). The bilingual nature of the tool
supports the increasing global relevance of Taiwan's capital market
disclosures, bridging language barriers and fostering international
investment confidence (Bender2021; Wen2023). Importantly, the
framework's scalable and replicable design facilitates adaptation to
other regional markets with similar linguistic or regulatory
complexities, while supporting continuous improvement through
incremental data updating.

Despite these advances, several limitations temper the generalizability
and direct applicability of our results. First, the training dataset,
though carefully curated and representative of various sectors and
languages, remains limited in size relative to the full scope of
Taiwan-listed companies' ESG disclosures. This constraint potentially
affects the LLM's exposure to rare or emerging ESG issues, which may
challenge comprehensive generalizability. Future research should explore
expanding the dataset longitudinally and diversifying sectoral coverage
to enhance robustness (GascoHernandez2022). Second, while the LLM
demonstrates improved accuracy, it is not immune to intrinsic model
biases or hallucination effects documented in large pretrained models
(Bender2021). Though fine-tuning and human review mitigate some risks,
ongoing vigilance and periodic recalibration are essential to guard
against propagating misleading or incomplete ESG interpretations. Third,
our study does not deeply investigate human-AI interaction dynamics,
including expert evaluators' acceptance, trust, and intuition
calibration when collaborating with LLM outputs. Kang2021 noted that
automation can unintentionally demotivate human participation or affect
judgment patterns; consequently, future work should analyze usability,
explainability, and effective interface integration to optimize human-AI
synergies in ESG evaluation.

Another consideration relates to domain specificity. ESG disclosures
often involve sector-specific jargon and evolving regulatory
requirements. While the Taiwan-specific fine-tuning addressed many
contextual nuances, rapid shifts in disclosure standards or emergent
social/environmental topics may necessitate frequent model updates or
adaptive learning mechanisms (Virmani2023; Manohar2023). In addition,
certain qualitative ESG factors such as corporate culture or governance
dynamics may remain challenging for AI to fully capture, underscoring
the continued importance of human interpretative insight. Thus, a hybrid
assessment paradigm, wherein LLM assistance accelerates routine analyses
while empowering experts to focus on high-complexity judgments, appears
most viable.

In conclusion, this study confirms that leveraging LLMs, properly
regionally adapted and fine-tuned, significantly enhances ESG disclosure
assessment accuracy, consistency, and efficiency for Taiwan-listed
companies. The results fill important empirical gaps in understanding
AI-assisted ESG evaluation in a multilingual, regulatory-specific
context, contributing novel technical and practical insights aligned
with prior literature (Bender2021; Kang2021; GascoHernandez2022;
Virmani2023). By addressing manual bottlenecks and standardization
issues, the proposed framework offers a scalable pathway toward improved
ESG transparency and accountability, ultimately supporting
sustainability governance both within Taiwan and comparable markets.
Future research should build on these foundations to explore deployment
experiences, refine human-AI collaboration, and expand across broader
geographic and sectoral boundaries to further advance the field.

\section{Conclusion}\label{conclusion}

This study presents the development and empirical validation of a novel
LLM-assisted framework for ESG disclosure assessment tailored
specifically to the bilingual and regulatory environment of
Taiwan-listed companies. Our findings confirm that the proposed approach
yields significant improvements in accuracy, consistency, and processing
efficiency compared to traditional manual evaluations, as evidenced by
statistically robust gains in accuracy (+7.4\%) and inter-rater
reliability (Fleiss' Kappa increase from 0.653 to 0.782), along with a
substantial reduction in assessment time by over 50\% (Table 2). These
results substantiate the testable hypotheses and underscore the critical
importance of regional customization and fine-tuning of large language
models to capture Taiwan's linguistic and regulatory nuances, addressing
key literature gaps identified by Bender et al.
(\citeproc{ref-Bender2021}{2021}) and Gasco-Hernandez, Gil-Garcia, and
Luna-Reyes (\citeproc{ref-GascoHernandez2022}{2022}).

Practically, this research offers a replicable and scalable tool that
can enhance the transparency and reliability of ESG disclosure
evaluation for regulators, investors, and rating agencies in Taiwan. It
also serves as a pioneering example of integrating advanced AI
techniques in sustainability governance within an Asian market context,
complementing human expertise without supplanting it, thereby mitigating
risks associated with human-AI collaboration highlighted by Kang et al.
(\citeproc{ref-Kang2021}{2021}).

Future research should investigate (1) longitudinal deployment and user
acceptance studies to understand the practical integration of LLM
assistance in real-world ESG workflows, (2) expansion of the framework
across other multilingual Asian markets to evaluate generalizability,
and (3) refinement of fairness and bias mitigation techniques to further
enhance model robustness and trustworthiness in sensitive ESG
evaluations. These directions will ensure continued advancement in
leveraging AI for sustainable business practices globally.

\section*{References}\label{references}
\addcontentsline{toc}{section}{References}

\phantomsection\label{refs}
\begin{CSLReferences}{1}{0}
\bibitem[\citeproctext]{ref-Attia2021}
Attia, Yasser A., Mohamed Samer, Mohamed A. Moselhy, Ahmed H. Arisha,
Ahmed A. Abdelqader, and Essam M. Abdelsalam. 2021. {``Influence of
Laser Photoactivated Graphitic Carbon Nitride Nanosheets and Nickel
Nanoparticles on Purple Non-Sulfur Bacteria for Biohydrogen Production
from Biomass.''} \emph{Journal of Cleaner Production} 299: 126898--98.
\url{https://doi.org/10.1016/j.jclepro.2021.126898}.

\bibitem[\citeproctext]{ref-Bender2021}
Bender, Emily M., Timnit Gebru, Angelina McMillan-Major, and Shmargaret
Shmitchell. 2021. {``On the Dangers of Stochastic Parrots.''}
\emph{Proceedings of the 2021 ACM Conference on Fairness,
Accountability, and Transparency}, 610--23.
\url{https://doi.org/10.1145/3442188.3445922}.

\bibitem[\citeproctext]{ref-FernndezEdreira2021}
Fernández-Edreira, Diego, Jose Liñares-Blanco, and Carlos
Fernandez-Lozano. 2021. {``Machine Learning Analysis of the Human Infant
Gut Microbiome Identifies Influential Species in Type 1 Diabetes.''}
\emph{Expert Systems with Applications} 185: 115648--48.
\url{https://doi.org/10.1016/j.eswa.2021.115648}.

\bibitem[\citeproctext]{ref-GascoHernandez2022}
Gasco-Hernandez, Mila, J. Ramon Gil-Garcia, and Luis F. Luna-Reyes.
2022. {``Unpacking the Role of Technology, Leadership, Governance and
Collaborative Capacities in Inter-Agency Collaborations.''}
\emph{Government Information Quarterly} 39: 101710.
\url{https://doi.org/10.1016/j.giq.2022.101710}.

\bibitem[\citeproctext]{ref-Kang2021}
Kang, Seonjun, Xiaojin (Jim) Liu, Yeongin Kim, and Victoria Yoon. 2021.
{``Can Bots Help Create Knowledge? The Effects of Bot Intervention in
Open Collaboration.''} \emph{Decision Support Systems} 148: 113601--1.
\url{https://doi.org/10.1016/j.dss.2021.113601}.

\bibitem[\citeproctext]{ref-Manohar2023}
Manohar, Sridhar, Justin Paul, Carolyn Strong, and Amit Mittal. 2023.
{``INNOSERV: Generalized Scale for Perceived Service Innovation.''}
\emph{Journal of Business Research} 160: 113723--23.
\url{https://doi.org/10.1016/j.jbusres.2023.113723}.

\bibitem[\citeproctext]{ref-Nawaz2021}
Nawaz, Saima, and Nasir Iqbal. 2021. {``How Cash Transfers Program
Affects Environmental Poverty Among Ultra-Poor? Insights from the BISP
in Pakistan.''} \emph{Energy Policy} 148: 111978--78.
\url{https://doi.org/10.1016/j.enpol.2020.111978}.

\bibitem[\citeproctext]{ref-OwusuAgyei2020}
Owusu-Agyei, Samuel, Godwin Okafor, Aruoriwo Marian Chijoke-Mgbame,
Paschal Ohalehi, and Fakhrul Hasan. 2020. {``Internet Adoption and
Financial Development in Sub-Saharan Africa.''} \emph{Technological
Forecasting and Social Change} 161: 120293--93.
\url{https://doi.org/10.1016/j.techfore.2020.120293}.

\bibitem[\citeproctext]{ref-Virmani2023}
Virmani, Naveen, Vaishali Agarwal, Koppiahraj Karuppiah, Satakshi
Agarwal, Rakesh D. Raut, and Sanjoy Kumar Paul. 2023. {``Mitigating
Barriers to Adopting Electric Vehicles in an Emerging Economy
Context.''} \emph{Journal of Cleaner Production} 414: 137557--57.
\url{https://doi.org/10.1016/j.jclepro.2023.137557}.

\bibitem[\citeproctext]{ref-Wen2023}
Wen, Shaojie, Wen-Chieh Cheng, Dongfeng Li, and Wenle Hu. 2023.
{``Evaluating Gas Breakthrough Pressure and Gas Permeability in a
Landfill Cover Layer for Mitigation of Hazardous Gas Emissions.''}
\emph{Journal of Environmental Management} 336: 117617--17.
\url{https://doi.org/10.1016/j.jenvman.2023.117617}.

\end{CSLReferences}




\end{document}
